% !TeX root = ../main.tex
\section{Conclusion}
\label{sec:conclusion}

We presented a first-order behavioral simulation framework for organic optoelectronic compute-in-memory inference that links materials-level characterization to task-level vision performance. By combining differential mapping, variability, retention, and peripheral constraints within a profile-driven workflow, the framework serves as a practical materials-to-system decision aid for evaluating deployment risk before full array hardware becomes available.

Several observations stand out. A 63-point iso-accuracy sweep reveals a two-phase constraint hierarchy: ADC resolution below 6 bits causes abrupt collapse ($S_{\text{ADC}}=0.98$ over the full grid), while within the operational envelope device-to-device variability becomes the binding factor ($S_{\text{D2D}}=0.92$). Fixed hardware-instance transfer, rather than nominal quantization, is the most prominent deployment risk in the present regime: standard HAT collapses on fresh arrays, whereas Ensemble HAT raises fresh-instance accuracy from 10.00~\% to $86.37 \pm 1.54\%$. In the canonical uniform-noise setting, quantization and nominal stochasticity are less limiting once scale recovery is active. By contrast, severe write non-linearity remains difficult to recover under the present gradient-scaling approximation ($27.72 \pm 0.82\%$), and determining whether physical organic devices enter this regime will require direct nonlinear-write measurements.

The framework translates reported device behavior into deployment-facing system risk through an explicit profile interface. Under a literature-anchored OPECT case study, zero-shot deployment reaches 88.53\% without any change to the evaluation workflow. The same structure should simplify future measured-device calibration while preserving model-level comparability within a common workflow.

\section*{Data Availability}
All datasets used in this study (CIFAR-10, CIFAR-100, and Flowers-102) are publicly available through standard PyTorch torchvision and HuggingFace datasets libraries. Source-data tables for the plotted figures, together with the submission-matched logs and inference traces, will be deposited in a reviewer-accessible repository at submission and released publicly at acceptance.

\section*{Code Availability}
The simulation framework, training and evaluation scripts, device-profile utilities, and plotting code will be provided to editors and reviewers through a reviewer-accessible archive at submission. A curated public release will be made available at \url{https://github.com/OrgOptEdge/compute_vit} upon acceptance.

\section*{Competing Interests}
The authors declare no competing interests.

\section*{Author Contributions}
S.L. conceived the framework, implemented the simulator, conducted all experiments, and wrote the manuscript.
